This paper characterizes rare-end inward localization in imbalanced ordinal
classification. A rare upper endpoint can be mislocalized toward interior
classes in a way that is not fully captured by exact-class accuracy or an
aggregate ordinal error. Routing, predictive location, and inward shrinkage
make this direction of failure visible and separate it from generic rare-class
underperformance.

The main replicated result is a direction-responsive classifier-head component.
Under the frozen A/C protocol, direction-only adaptation reduced rare-end
localization error in all four confirmatory backbone seeds in each of the
RetinaMNIST CE, RetinaMNIST RPS, Solar CE, and Solar RPS settings. Because the
intervention holds the representation, original class-specific row norms, and
biases fixed, this result provides evidence that head direction is an actionable
component of rare-end localization in the studied settings. It does not define
a universally preferable classifier or imply a global predictive, calibration,
or uncertainty-quantification improvement.

The associated geometry analysis gives a narrower conclusion. Endpoint-specific
representation geometry added explanatory value strongly in Solar, was partial
in RetinaMNIST CE, and was not supported in RetinaMNIST RPS. The geometry
diagnostic is therefore setting-dependent rather than a universal explanation
for direction-only recovery. The controlled RetinaMNIST factor study similarly
identified balanced sampling as the dominant observed adaptation signal only in
its frozen one-backbone setting. The rare-support severity study showed weaker
direct class-4 evidence at lower support without a clean monotonic localization
dose response.

Taken together, the results support an empirical mechanism account rather than
a deployment method: rare-end inward localization has a reproducible
direction-responsive head component, while representation diagnostics and
support effects require setting-specific interpretation. Future work can test
this account on broader ordinal datasets and backbones, develop richer
representation diagnostics, and determine when a head-level response remains
useful under different imbalance regimes.
